\documentclass[12pt]{article}
\usepackage[utf8]{inputenc}
\usepackage[T1]{fontenc}
\usepackage{hyperref}
\title{Tamagario Language Resources for K-12 Curriculum Development}
\author{Compiled by Claude Sonnet 5 (claude-sonnet-5) for Andr\'as Kornai}
\date{}
\begin{document}
\maketitle

\section{Introduction}
Tamagario (ISO 639-3 \texttt{tcg}) is a Trans-New Guinea language of Mappi Regency, South Papua, Indonesia. Wikipedia gives about 3,500 speakers, citing a 1987 figure. Joshua Project gives population estimates of about 6,000. Sources disagree on its relation to Kayagar (\texttt{kyt}), which has its own ISO code; Joshua Project lists Tamario and Wiyagar as Kayagar dialects. Yogo is treated as a dialect by some sources and as a separate language by others. Speakers live in swampy lowlands along the Digul and Mapi Rivers. This survey found very little material, and nothing designed for schools.

\section{Primary Educational Resources}
No primers, textbooks or school programmes for Tamagario turned up in the searches.

\section{Linguistic and Reference Materials}
\begin{itemize}
\item Wikipedia, Tamagario language (overview and classification): \url{https://en.wikipedia.org/wiki/Tamagario_language}
\item Wikipedia, Kayagar languages (includes a comparative word list with Tamagario): \url{https://en.wikipedia.org/wiki/Kayagar_languages}
\item New Guinea World, Tamagario (survey entry with classification and dialect notes): \url{https://newguineaworld.linguistik.uzh.ch/families/trans-new-guinea/cook-river-kolopom/cook-river/gondu-river/central-gondu-river/tamagario}
\item ASJP, Tamagario word list: \url{https://asjp.clld.org/languages/TAMAGARIO_2}
\end{itemize}

\section{Religious and Community Resources}
\begin{itemize}
\item Joshua Project, Tamagario (Buru) in Indonesia (people-group profile, not a language resource as such): \url{https://joshuaproject.net/people_groups/15222}
\item Joshua Project, Kayagar in Indonesia (related group): \url{https://joshuaproject.net/people_groups/12595}
\end{itemize}

\section{Digital and Audio Resources}
No audio, video or software resources turned up in the searches.

\section{Curriculum Development Context}
Little published material exists. A curriculum would have to start from new documentation: orthography development, word lists and community-recorded texts. Anyone doing this should first settle how Tamagario relates to Kayagar and Yogo.
\end{document}